% Generated draft: requires mathematical and source review.
\documentclass[11pt]{article}
\usepackage[margin=25mm]{geometry}
\usepackage{fontspec}
\setmainfont{Latin Modern Roman}
\usepackage{amsmath,amssymb,mathtools,unicode-math}
\setmathfont{Latin Modern Math}
\usepackage{enumitem,needspace,xurl,hyperref}
\setlength{\parindent}{0pt}
\setlength{\parskip}{5pt}
\setlength{\emergencystretch}{3em}
\newcommand{\N}{\mathbb{N}}
\newcommand{\Z}{\mathbb{Z}}
\newcommand{\Q}{\mathbb{Q}}
\newcommand{\R}{\mathbb{R}}
\newcommand{\C}{\mathbb{C}}
\newcommand{\F}{\mathbb{F}}
\newcommand{\A}{\mathbb{A}}
\newcommand{\PP}{\mathbb{P}}
\newcommand{\E}{\mathbb{E}}
\newcommand{\eps}{\varepsilon}
\newcommand{\dd}{\,\mathrm{d}}
\begin{document}
\section*{UnsolvedMath ID: 30002066}
\textbf{Source problem code:} catalogue entry
\subsection{Definitions and mathematical statement}
The supplied record contains two groups of questions whose full formulation requires additional conventions.

\begin{enumerate}
\item A closed \(3\)-manifold is a compact \(3\)-manifold without boundary. After fixing a notion of triangulation, an encoding, and an equivalence relation on triangulations, let \(\mathcal T_N\) denote the resulting finite sampling space of triangulations of closed \(3\)-manifolds having exactly \(N\) tetrahedra. Let
\[
\mathcal S_N=\{T\in\mathcal T_N:|T|\cong S^3\},
\]
where \(|T|\) denotes the underlying topological space and \(\cong\) denotes homeomorphism.

The sampling question asks whether there is a randomized algorithm that, for every positive integer \(N\) with \(\mathcal T_N\ne\varnothing\), outputs an element of \(\mathcal T_N\) in time polynomial in \(N\), with
\[
\Pr(A(N)=T)=\frac{1}{|\mathcal T_N|}
\qquad\text{for every }T\in\mathcal T_N.
\]
It also asks the corresponding question with \(\mathcal T_N\) replaced by \(\mathcal S_N\), allowing a separate algorithm. The source asks for uniform sampling and supplies no approximation tolerance.

The source does not specify whether triangulations are simplicial complexes or allow tetrahedral face identifications, whether they are labelled or considered up to isomorphism, or whether connectedness or orientability is required. It also does not specify the encoding, the running-time guarantee, or the treatment of values of \(N\) for which a sampling space is empty. Thus the displayed conditions express the sampling target conditional on conventions that remain unresolved.

\item A random graph on a fixed labelled vertex set is exchangeable if its distribution is invariant under relabelling the vertices; for a countably infinite vertex set, this means invariance under every permutation moving only finitely many vertices. For a graph \(G\), its clique complex \(\mathrm{Cl}(G)\) has as simplices the finite vertex subsets whose distinct vertices are pairwise adjacent. For a basepoint \(x\in|\mathrm{Cl}(G)|\), write
\[
\pi_k\bigl(|\mathrm{Cl}(G)|,x\bigr),\qquad k\geq 1,
\]
for its homotopy groups.

The source asks how these homotopy groups behave for exchangeable random graphs and whether the resulting random objects satisfy a universality property. However, it specifies neither a graph distribution nor a finite or infinite vertex regime, an asymptotic parameter, a basepoint or connected component, a probabilistic mode of conclusion, or a definition of universality. These questions therefore remain requests for a more precise formulation; the supplied evidence does not determine a single quantified homotopy-group or universality assertion.
\end{enumerate}
\subsection{Short English statement}
The record asks for polynomial-time exact uniform sampling of closed \(3\)-manifold triangulations with a prescribed number of tetrahedra, including the restriction to \(S^3\). It separately asks about homotopy groups and universality for clique complexes of exchangeable random graphs. Sampling conventions, the random graph law, and the meaning of universality are unspecified. The supplied open status and source-verification claims require independent review.
\subsection{Sources}
\begin{itemize}
\item Supplied local UnsolvedMath registry record, permanent ID \(30002066\), source problem code \texttt{catalogue entry}, title: Random Triangulations and Exchangeable Clique Complexes. Collection: Oberwolfach Reports Open Problems. No individual proposer is supplied; the proposed year is recorded as 2012.
\item Supplied report citation: \emph{Triangulations (2013)}, Oberwolfach Reports, DOI 10.4171/owr/2012/24, pp.\ 1405--1486. \url{https://doi.org/10.4171/owr/2012/24}. The background attributes the questions to workshop open-problem discussions documented in Oberwolfach Reports.
\item Supplied literature source labelled ``original OWR report'': \url{https://doi.org/10.4171/owr/2012/24}.
\item Supplied literature source labelled ``random-clique-complex literature'': \url{https://doi.org/10.1016/j.disc.2008.11.012}. The record supplies no author, title, or specific supporting result for this reference.
\item Literature assessment dated exactly \texttt{2026-08-21}: the supplied curation labels the problem open and reports that no general resolution was found for the combined sampling, \(3\)-sphere, homotopy-group, and universality questions. This is supplied dated curation, not an independently verified literature assessment. The background expressly requests independent verification of time-sensitive status and cited claims.
\item The registry labels the statement \texttt{corrected\_verified} and asserts that it was checked against the official Oberwolfach report. These are supplied provenance claims; no independent source visit or mathematical verification is asserted here.
\end{itemize}
\end{document}
